\chapter{Introduction}
\label{ch:introduction}

\section{Motivation and Context}
\label{sec:motivation}

Artificial intelligence has spent the last several years moving away from uni-modal processing and toward representations that span multiple modalities at once. Foundational vision-language models, \gls{clip} chief among them, showed that projecting disparate modalities --- two-dimensional images and natural language text, in that case --- into a single, continuous latent space produces representations that are useful for far more than either modality alone. That shared semantic space is what makes zero-shot classification, cross-modal retrieval, and text-guided visual synthesis possible in the first place, bridging symbolic linguistic concepts and raw visual input.

Three-dimensional representation, however, has not progressed at a comparable pace. Multimodal alignment is fairly mature in 2D vision, but extending the same ideas into 3D --- point clouds, polygonal meshes, volumetric voxel grids --- remains an open problem, and one that matters: robotics, autonomous driving, \gls{arvr}, and embodied AI all depend on it. An intelligent agent trying to act in the physical world needs to connect a phrase like \emph{``a wooden dining chair with curved legs''} to the actual geometric structure it refers to, not just recognise the words.

Aligning 3D representations with natural language is substantially more difficult than the 2D case, for two fairly distinct reasons:
\begin{itemize}
    \item \textbf{Geometric Dissimilarity:} 2D images are regular pixel grids, which line up naturally with sequence-based text tokens. 3D representations do not offer anything like that regularity --- they are non-Euclidean, sparse, permutation-invariant, and topologically irregular.
    \item \textbf{Data Scarcity:} 2D vision-language models are trained on billions of scraped image-text pairs. Nothing close to that scale exists for paired 3D-text data, which rules out joint end-to-end multimodal pre-training from scratch as a practical option.
\end{itemize}

One way around the need for massive joint pre-training is \emph{a posteriori} cross-modal alignment: instead of training multimodal architectures end-to-end, this line of work asks whether independently trained, uni-modal 3D encoders (PointNet, SparseConv) and pre-trained language models (OpenCLIP, \gls{roberta}, \gls{bert}) already build topological structures that happen to share a common geometry. \textcite{hadgi2025escaping}, in \emph{``Escaping Plato's Cave''}, argued that if independently trained uni-modal spaces do capture shared real-world concepts, then post-hoc mathematical projections --- \gls{cca} and affine mappings, specifically --- should be able to align them, with no joint retraining required.

This thesis is built around the following question: what are the mechanics, the capacity, and the fundamental limits of post-hoc 3D-text alignment? Comparing how well uni-modal spaces can be aligned --- through supervised linear projections on one hand, and unsupervised relational frameworks such as \gls{ot} on the other --- should reveal whether 3D shapes and natural language actually share an intrinsic semantic geometry, or whether any apparent alignment is an artefact of supervision.

\section{Problem Statement}
\label{sec:problem_statement}

Post-hoc multimodal alignment has advanced considerably, but a reliable way of mapping uni-modal 3D representations onto natural language remains unsolved. Standard vision-language models circumvent this limitation by training jointly and contrastively on massive paired datasets from the start; post-hoc alignment instead attempts to map latent spaces that were trained completely independently, without modifying the underlying backbone encoders at all. Even with some pioneering work already in this space, two fundamental questions remain open:

\begin{enumerate}

\item \textbf{The Curse of Paired Supervision:} The framework proposed by \textcite{hadgi2025escaping} shows that supervised linear projections --- \gls{cca} followed by an affine transformation --- can extract shared structure between 3D and text representations. However, their method requires thousands of explicitly paired anchors to do so, which raises the following question: how much of that alignment reflects something intrinsic to the independently learned latent spaces, and how much is simply forced into existence by the supervision itself?

\item \textbf{The Effectiveness of Relational Unsupervised Alignment:} \gls{ot}, in its \gls{gw} formulation, is designed to align two metric spaces using nothing but their intra-domain relational structure --- no pairwise supervision at all, in principle. \gls{gw} has demonstrated success in graph matching and 2D vision-language alignment, but it remains untested whether that success carries over to high-dimensional, complex 3D-text latent spaces.

\end{enumerate}

This thesis directly addresses both of these questions, exploring the empirical and theoretically-grounded limits of supervised linear projections and unsupervised relational optimal transport for 3D-text cross-modal alignment.

\subsection{Research Questions}

\label{subsec:research_questions}

Three research questions run through the rest of this thesis:

\begin{itemize}

\item \textbf{RQ1 (Supervised Baseline Validation):} How well do supervised linear projections (\gls{cca} with affine transformations) of a 3D point cloud embedding space align with various natural language latent spaces --- \gls{clip}, \gls{roberta}, \gls{bert}?

\item \textbf{RQ2 (Feasibility of Unsupervised Optimal Transport):} Can unsupervised or weakly-supervised Gromov-Wasserstein-based methods (\gls{sgw}, \gls{risgw}, \gls{ssrisgw}, \gls{fgw}, \gls{lrgw}) learn a valid alignment between 3D and text latent spaces without any explicit anchor pairs, or with only minimal anchor-based supervision?

\item \textbf{RQ3 (Structural and Topological Analysis):} If unsupervised relational optimal transport does fail to achieve meaningful cross-modal alignment, what mathematical properties of the latent spaces --- or of the transport mechanisms themselves --- are responsible?

\end{itemize}

\subsection{Thesis Objectives and Scope}

\label{subsec:thesis_scope}

To address these questions, this thesis pursues four concrete objectives:

\begin{itemize}

\item \textbf{Implement and Reproduce the Supervised Baseline:} Build an end-to-end evaluation framework on a large paired 3D-object dataset (Objaverse, with Cap3D captions) to reproduce and benchmark \textcite{hadgi2025escaping}'s supervised linear mapping across multiple 3D and text encoder configurations.

\item \textbf{Develop and Deploy Unsupervised OT Algorithms:} Implement and evaluate five Gromov-Wasserstein variants (\gls{sgw}, \gls{risgw}, \gls{ssrisgw}, \gls{fgw} and \gls{lrgw}), combined with the scale-matching normalisations proposed by \textcite{li2026scale}.

\item \textbf{Conduct Exhaustive Ablation Studies:} Systematically probe hyperparameter ranges and feature-space effects --- with and without projections, subspace dimensionality, regularisation strength --- to isolate the sources of computational and performance bottlenecks.

\item \textbf{Explain the Observed Behaviour via Existing Theory:} Develop grounded explanations for why unsupervised alignment succeeds or fails, through a detailed analysis of domain-specific metric distribution properties, \gls{cka} structural similarity, and transport plan entropy, drawing on established results from the optimal transport literature.

\end{itemize}

\section{Contributions}
\label{sec:contributions}

This thesis makes the following technical, empirical, and methodological contributions to cross-modal representation learning and 3D-text latent alignment.

\subsection{Key Contributions}

\label{subsec:contributions}

\begin{itemize}

\item \textbf{Faithful Replication and Benchmarking of Supervised Linear Alignment:} We build an end-to-end feature extraction and alignment pipeline covering $46{,}675$ paired 3D objects from Objaverse with Cap3D captions, faithfully replicating \textcite{hadgi2025escaping}'s \gls{svd}-based \gls{cca} and affine alignment baseline across six encoder combinations (PointNet and SparseConv, each paired with OpenCLIP, \gls{roberta}, and \gls{bert}). This establishes a solid evaluation benchmark, peaking at $35.7\%$ Top-5 retrieval for PointNet $\times$ OpenCLIP.

\item \textbf{Exhaustive Evaluation of Optimal Transport Frameworks:} We implement and benchmark five \gls{gw}-based \gls{ot} variants --- \gls{sgw}, \gls{risgw}, \gls{ssrisgw}, \gls{fgw}, and \gls{lrgw} via OTT-JAX --- incorporating \textcite{li2026scale}'s scale-matching normalisation. An exhaustive grid search plus thirteen independent ablation experiments together demonstrate, systematically rather than anecdotally, that these methods fail to achieve meaningful alignment: the best fully unsupervised variant (\gls{sgw}) reaches only $3.2\%$ Top-5, and even the best semi-supervised variant (\gls{fgw}, using a limited anchor set) reaches only $4.4\%$ --- both nowhere near sufficient to align 3D and text latent spaces, and both far below the fully supervised \gls{cca}+Affine baseline of $35.7\%$.

\item \textbf{Exploration of Non-Linear Supervised Alignment Extensions:}
As an original contribution beyond linear mapping, we implement and evaluate
two supervised non-linear alignment frameworks: \gls{kcca} with
Nystr\"{o}m kernel matrix approximation, extending linear \gls{cca} into
non-linear territory via an \gls{rbf} kernel; and \gls{rml}, which learns a
Mahalanobis distance metric by gradient descent on the Cholesky factor of a
\gls{spd} matrix --- a more flexible, non-linear alternative to directly mapping
feature spaces.

\item \textbf{Empirical Demonstration of Structural Non-Isometry:} By analysing transport plan entropy, mass concentration ratios, self-\gls{gw} cost ratios, and unaligned \gls{cka} structural similarity scores (Hadgi et al.'s Figure 2 reports $\approx 0.12$ for PointNet$\times$CLIP, against $0.31$ for the jointly-trained OpenShape$\times$CLIP pair \cite{hadgi2025escaping}), we develop a grounded explanation, drawing on established Gromov-Wasserstein theory, for why unsupervised relational optimal transport fails in this setting. Uni-modal 3D point cloud latent spaces and text spaces fundamentally violate the assumption of \emph{structural isometry}, which is what causes quadratic optimal transport solvers to become trapped in non-informative local minima.

\end{itemize}

\section{Thesis Outline}
\label{subsec:thesis_structure}


The rest of this thesis is organised as follows:

\begin{itemize}
    \item \textbf{Chapter \ref{ch:background} (Background and Related Work):}
    Works through the core theoretical foundations --- 3D representation learning,
    multimodal vision-language models, Canonical Correlation Analysis, and
    Optimal Transport theory, from Monge-Kantorovich formulations and
    Wasserstein metrics through to Gromov-Wasserstein relational alignment.

    \item \textbf{Chapter \ref{ch:experimental_setup} (Experimental Setup):}
    Covers the experimental methodology: dataset preparation (Objaverse and
    Cap3D), backbone encoder architectures, the high-performance computing cluster
    setup, and the cross-modal retrieval evaluation metrics used throughout.

    \item \textbf{Chapter \ref{ch:baseline} (Supervised Linear Baseline):}
    Details the implementation, replication, and performance of the supervised
    linear alignment pipeline (\gls{cca} and affine transformations) across all six
    encoder combinations.

    \item \textbf{Chapter \ref{ch:nonlinear} (Non-Linear Alignment Methods):}
    Presents the two supervised non-linear extensions --- Kernel \gls{cca} with
    Nystr\"{o}m approximation and Riemannian Metric Learning --- evaluated at
    full dataset scale and compared against the linear baseline.

    \item \textbf{Chapter \ref{ch:gw_alignment} (Gromov-Wasserstein Alignment
    and Structural Analysis):} Investigates unsupervised relational Optimal
    Transport formulations, walks through thirteen diagnostic ablation experiments
    across \gls{gw} variants, and lays out an empirically-grounded non-isometry
    analysis explaining the fundamental limits of relational \gls{ot}.

    \item \textbf{Chapter \ref{ch:conclusions} (Conclusion and Future Work):}
    Pulls together the thesis's main findings, discusses what they mean for 3D
    multimodal alignment more broadly, and outlines directions worth pursuing
    next.
\end{itemize}